\begin{figure*}[!ht]
\centering
\includegraphics[width=\linewidth]{./Figure/Phase.png}
\caption{
    (a--c) Phase diagrams of L--S TT copolymers in the $N\chi$--$\bar{\phi}_{\mathrm A}$ plane for $\phi_{\mathrm a}=0.25$, $0.50$, and $0.75$, respectively.
    Red lines denote SCFT results for the L--S TT system, gray lines denote the reference phase diagram of linear AB diblock copolymers, and blue dashed lines denote the RPA spinodal boundaries~\cite{Zhao2024Phase}.
    Insets show representative three-dimensional component-density isosurfaces of the main ordered phases, including FCC, BCC, C, G, $\mathrm{O}^{70}$, and L.
}
\label{fig:PhaseDiagram}
\end{figure*}



Sliding reshapes the phase boundaries of L--S TT copolymers.

Fig.~\ref{fig:PhaseDiagram} presents the SCFT phase diagrams of L--S TT copolymers for representative sliding side-chain fractions, $\phi_{\mathrm a}=0.25,\,0.50,\,0.75$, together with the spinodal boundaries (blue dashed lines) and SCFT ODT boundaries, and compares them with the phase diagram of conventional linear AB diblock copolymers (gray lines).
Because L--S TT molecules undergo a transition from star topology locking to linear topology locking as $\phi_{\mathrm a}$ changes, the spinodal lines at $\phi_{\mathrm a}=0.25$ and $\phi_{\mathrm a}=0.50$ lie slightly below those of linear AB diblock copolymers, whereas at $\phi_{\mathrm a}=0.75$ they almost coincide.

The sliding ensemble also modifies the conformational resistance of the homogeneous phase to composition fluctuations, as reflected by $\Delta S_{\mathrm{conf}}$ in the preceding analysis. 
As a result, for all values of $\phi_{\mathrm a}$, the SCFT ODT boundaries of L--S TT copolymers lie significantly above the corresponding SCFT ODT boundaries of linear AB diblock copolymers and also above the RPA spinodal lines of the L--S TT system.


% 图~\ref{fig:PhaseDiagram} 的另一重要特征是，当 $\phi_\mathrm{a}\neq 0$ 时，相图结构呈现出显著的非对称性，各相区不再关于 $\bar{\phi}_\mathrm{A}=0.5$ 对称。  
% 这种不对称性源于滑动侧链的引入打破了传统 AB 二嵌段共聚物的拓扑对称性，使体系在结构上类似于引入具有 A$_2$B 型特征的不对称链段，从而诱导出宏观相图的偏斜。  
% 在这种情况下，$\bar{\phi}_\mathrm{A}<0.5$ 与 $\bar{\phi}_\mathrm{A}>0.5$ 分别对应滑动嵌段主要位于胶束核心（core）与壳层（shell）时的相行为。  
% 因此，相图在低组分一侧（$\bar{\phi}_\mathrm{A}<0.5$）整体收缩，而在高组分一侧（$\bar{\phi}_\mathrm{A}>0.5$）则相应扩展，这一特征在高 $N\chi$ 区域尤为明显，例如高 $\phi_\mathrm{A}$ 区域的 BCC 相区。
Another important feature of Fig.~\ref{fig:PhaseDiagram} is that, when $\phi_{\mathrm a}\neq 0$, the phase boundaries are no longer symmetric about $\bar{\phi}_{\mathrm A}=0.5$.
This asymmetry originates from the sliding side chain, which breaks the topological symmetry of AB diblock copolymers.
As a result, the sliding side chain introduces an A$_2$B-like architectural asymmetry, which leads to a skewed phase diagram at the macroscopic level.
Accordingly, relative to the linear AB diblock reference, the phase diagram is compressed for $\bar{\phi}_{\mathrm A}<0.5$ and expanded for $\bar{\phi}_{\mathrm A}>0.5$.
This asymmetry reflects two different packing environments of the sliding A segment: it is located predominantly in the micellar core for $\bar{\phi}_{\mathrm A}<0.5$, whereas it mainly contributes to the shell for $\bar{\phi}_{\mathrm A}>0.5$.
This feature is particularly evident in the high-$N\chi$ region, for example in the BCC phase region at high $\bar{\phi}_{\mathrm A}$.


\begin{figure*}[!ht]
    \centering
    \includegraphics[width=\linewidth]{./Figure/Gyroid.png}
    \caption{
        (a) Three-dimensional isosurfaces of the A- and B-component density distributions at $\phi_{\mathrm a}=0.50$, $\bar{\phi}_{\mathrm A}=0.30$, and $N\chi=30$.
        (b) Spatial distributions of the blocks in L--S TT copolymers for $\phi_{\mathrm a}=0.25$, $0.50$, and $0.75$.
        Rows show two-dimensional density slices along the main diagonal ($x=y$), corresponding to the light-gray section in panel (a), for A1, A2, A3, and B, respectively.
        The densities of all A blocks are normalized by their corresponding overall volume fractions, $\phi_{\mathrm a}$ or $(1-\phi_{\mathrm a})$.
        (c) Probability distributions of the sliding degree of freedom, $P_\Lambda(f)$, for different $\phi_{\mathrm a}$ values, as indicated by color.
    }
    \label{fig:Gyroid}
\end{figure*}


% 为进一步揭示滑动拓扑自由度如何在分子尺度上响应上述相图变化，我们在 Gyroid 结构中系统分析各链段及拓扑参数 $f$ 的空间分布，从而刻画滑动侧链在三维有序结构中的空间自适应调节机制。
% 如图~\ref{fig:Gyroid} 所示，在设定 Ia$\overline{3}$d 空间群对称性下，在确定的组分比 $\bar{\phi}_\mathrm{A} = 0.30$，相互作用参数 $N\chi = 30.0$ 的条件下，对比 $\phi_{\mathrm{a}} = 0.25$, $0.50$, $0.75$ 的三种可滑动侧链体积分数的 L-S TT，均形成由可滑动链构成胶束核心的 Gyroid 相结构。
% 图~\ref{fig:Gyroid}a 展示了以 $\phi_{\mathrm{a}}=0.50$ 为例的典型形貌的三维视图。
% 图~\ref{fig:Gyroid}b 展示了在 $\phi_\mathrm{a}=0.25, 0.50, 0.75$ 三种不同侧链体积分数下，L–S TT 共聚物在 $(110)$ 晶面上的二维截面（见图~\ref{fig:Gyroid}a 中的灰色观测面）上的各嵌段 $m$ 的密度分布 $\phi_m(\bm r)$。
% 这些链段密度分布可由公式~(\ref{eq:BlockDensity}) 计算得到。
% 为便于不同 $\phi_\mathrm{a}$ 条件之间的对比，各 A 组分链段的密度已进行标准化处理，分别表示为 $\phi_\mathrm{A1}/\phi_\mathrm{a}$、$\phi_\mathrm{A2}/(1-\phi_\mathrm{a})$、$\phi_\mathrm{A3}/(1-\phi_\mathrm{a})$。  
% 图中每列对应不同的 $\phi_\mathrm{a}$，每行依次表示滑动侧链 A1、主链中的 tail 链段 A2、trunk 链段 A3 以及主链上的 B 组分的空间分布形态。
% 一般地，滑动侧链 $\phi_\mathrm{A1}$ 主要集中在 A 域的核心位置，形成致密的 A 核心区；trunk 链段靠近分子的接枝点，$\phi_\mathrm{A3}$ 因而位于 A/B 界面附近，包覆在 A 核外层。
% 此外，在不同滑动比例下 $\phi_\mathrm{B}$ 的空间分布形态基本一致，始终填充于 A 相胶束之间的间隙，表明不同 $\phi_\mathrm{a}$ 条件下胶束尺寸基本相同。
% 因此，不同 $\phi_\mathrm{a}$ 的差异主要体现在胶束内部各链段的分布上：
% 如 $\phi_\mathrm{a} = 0.25$ 时，$\phi_{\mathrm{A1}}$ 与 $\phi_{\mathrm{A2}}$都集中在 A 域的核心区域，对应于 RPA 预测的 “Y” 型 $\mathrm{A}_2\mathrm{AB}$ 构象，其中 tail 嵌段长度与滑动侧链长度相当\cite{Zhao2024Phase}。
% 而随 $\phi_\mathrm{a}$ 的增大，$\phi_\mathrm{A2}$ 的空间分布密度峰值逐渐由 A 核心区域向 A/B 界面外移。
% 因此在$\phi_\mathrm{a} = 0.75$的密度热力图中，$\phi_\mathrm{A1}$、$\phi_\mathrm{A2}$ 和 $\phi_\mathrm{A3}$ 呈现出由内至外依次嵌套的空间分布。
% 这种变化趋势表明，随着 $\phi_\mathrm{a}$ 的增大，L–S TT 分子趋向于形成更接近 AB 线型 diblock 的构型，以降低整体自由能。
% 对应的，$\phi_\mathrm{A1}$的空间分布则越来越向 A 组分的中间区域聚集；
% 而与 B 组分直接连接的 $\phi_\mathrm{A3}$ 始终分布于 A 组分区域的外壳层，且其空间位置基本不随 $\phi_\mathrm{a}$ 变化而改变，显示出其分布的相对稳定性。
To resolve the molecular origin of these phase diagram changes, we analyze the spatial distributions of the four blocks $m$ and the topology distribution $P_\Lambda(f)$ in the Gyroid phase.
This allows us to characterize the spatially adaptive regulation of the sliding side chain.
As shown in Fig.~\ref{fig:Gyroid}, under the prescribed Ia$\overline{3}$d space-group symmetry, L--S TT copolymers with three representative sliding side-chain fractions, $\phi_{\mathrm a}=0.25$, $0.50$, and $0.75$, all form a Gyroid phase with the micellar core composed of the sliding A component at $\bar{\phi}_{\mathrm A}=0.30$ and $N\chi=30.0$.
Fig.~\ref{fig:Gyroid}a shows a representative three-dimensional morphology for $\phi_{\mathrm a}=0.50$.
Fig.~\ref{fig:Gyroid}b shows two-dimensional density slices on the $(110)$ plane, corresponding to the gray observation plane in Fig.~\ref{fig:Gyroid}a, for the density distributions $\phi_m(\bm r)$ of each block $m$ at $\phi_{\mathrm a}=0.25$, $0.50$, and $0.75$.
These block density distributions are calculated from Eq.~(\ref{eq:BlockDensity}).
For comparison of the spatial block-density distributions at different $\phi_{\mathrm a}$, we rescale $\phi_{\mathrm{A1}}$ by $\phi_{\mathrm a}$ and both $\phi_{\mathrm{A2}}$ and $\phi_{\mathrm{A3}}$ by $1-\phi_{\mathrm a}$, corresponding to the prescribed fractions of A segments in the side-chain block and the main-chain A block, respectively.
This reduces composition driven differences in the overall density levels while preserving the spatial morphology and the relative distribution of A2 and A3.

In Fig.~\ref{fig:Gyroid}b, each column corresponds to a different $\phi_{\mathrm a}$, and the rows show, from top to bottom, the spatial distributions of the sliding side chain A1, the tail block A2, the trunk block A3, and the B block on the main chain.
In general, the sliding side chain $\phi_{\mathrm{A1}}$ is mainly concentrated in the core of the A domain, forming a dense A-rich core.
Because the trunk block is adjacent to the grafting point, $\phi_{\mathrm{A3}}$ is located near the A/B interface and forms an outer layer around the A core.
In addition, the spatial distribution of $\phi_{\mathrm B}$ remains nearly unchanged for different $\phi_{\mathrm a}$ and always fills the interstitial regions between A-rich micelles, indicating that the micellar size is almost unchanged as $\phi_{\mathrm a}$ varies.
Therefore, the main difference among different $\phi_{\mathrm a}$ lies in the internal partitioning of the A blocks within the micelles.
For $\phi_{\mathrm a}=0.25$, both $\phi_{\mathrm{A1}}$ and $\phi_{\mathrm{A2}}$ are concentrated in the core region of the A domain, corresponding to the ``Y''-shaped $\mathrm{A}_2\mathrm{AB}$ conformation predicted by RPA, where the tail block has a length comparable to that of the sliding side chain~\cite{Zhao2024Phase}.
As $\phi_{\mathrm a}$ increases, the density maximum of $\phi_{\mathrm{A2}}$ gradually shifts outward from the A-core region toward the A/B interface.
Consequently, for $\phi_{\mathrm a}=0.75$, the density maps show a nested spatial arrangement of $\phi_{\mathrm{A1}}$, $\phi_{\mathrm{A2}}$, and $\phi_{\mathrm{A3}}$ from the inner core to the outer interface.
This trend indicates that, with increasing $\phi_{\mathrm a}$, L--S TT molecules tend to adopt configurations closer to the linear AB diblock limit in order to lower the total free energy.
Correspondingly, the spatial distribution of $\phi_{\mathrm{A1}}$ becomes increasingly localized toward the central region of the A domain.
By contrast, $\phi_{\mathrm{A3}}$, which is directly connected to the B block, remains located in the outer shell of the A domain and shows little change in spatial position with varying $\phi_{\mathrm a}$.


% 对比不同 $\phi_\mathrm{a}$ 下 Gyroid 相的全空间滑动分布 $P_\Lambda(f)$ 的变化情况，可以进一步表征体系中滑动自由度带来的拓扑结构自适应特征。
% 固定  $N\chi = 30.0$，如图~\ref{fig:Gyroid}c，发现不同 $\phi_{\mathrm a}$ 条件下的 $P_{\Lambda}(f)$ 均呈现唯一的极大值，表明各拓扑 $f_j$ 并非等概率出现，而体系对某一特定拓扑具有明显偏好。
% 该峰值对应的拓扑定义为最优拓扑 $f^{*}$，并在图~\ref{fig:Gyroid}c 中以虚线标示。
% 随着可滑动链段比例 $\phi_\mathrm{a}$ 的增加，最优拓扑 $f_\Lambda^\ast$ 的取值逐渐降低。
% 这表明当可滑动侧链变长时，L-S TT 自适应形成类似 AB 二嵌段共聚物的线性拓扑结构，而非保持滑动链的 $\mathrm{A_2 AB}$ 星型拓扑结构，我们称其为线型拓扑择优。
% 这是 L-S TT 在自组装过程中，为了兼顾界面曲率匹配与填充挫折的缓解，
% 而自发调整滑动拓扑所形成的选择结果，与无规相近似理论（RPA）中的预测相符\cite{Zhao2024Phase}。
The variation of the full-space sliding distribution $P_\Lambda(f)$ for the Gyroid phase at different $\phi_{\mathrm a}$ further characterizes the adaptive topology selection induced by the sliding degree of freedom.
At fixed $N\chi=30.0$, as shown in Fig.~\ref{fig:Gyroid}c, $P_\Lambda(f)$ exhibits a single maximum for each $\phi_{\mathrm a}$, indicating that different topologies $f_j$ are not equally sampled and that the system has a preference for a specific topology.
The topology corresponding to this maximum is defined as the preferred topology $f_\Lambda^\ast$ and is marked by the dashed lines in Fig.~\ref{fig:Gyroid}c.
As the sliding chain fraction $\phi_{\mathrm a}$ increases, the value of $f_\Lambda^\ast$ gradually decreases.
This indicates that, when the sliding side chain becomes longer, L--S TT copolymers preferentially adopt a more linear AB-diblock-like topology rather than maintaining an $\mathrm{A_2AB}$ star-like topology, which we refer to as linear topology preference.
This preference results from the spontaneous adjustment of sliding topology during self-assembly, allowing the system to respond to variations in interfacial curvature and heterogeneous packing constraints, and is consistent with the prediction of the RPA~\cite{Zhao2024Phase}.
